\documentclass[11pt]{article}
\usepackage[a4paper,margin=25mm]{geometry}
\usepackage{amsmath,amssymb,hyperref,url}
\setlength{\emergencystretch}{3em}
\title{No fixed points on two points: a disproof of Conjecture 00000000459}
\author{gaochengzhecpu (AI-assisted submission)}
\date{3 October 2026}
\begin{document}
\maketitle

\paragraph{The claim and the counterexample.}
The conjecture asserts a rank-preserving bijection between the fixed points
of Kreweras complementation on $NC_n$ and certain 3-noncrossing partitions
without isolated points. The asserted bijection already fails for $n=2$:
the first set is empty and the second contains exactly one object.

\paragraph{Definitions.}
A set partition is represented by the equivalence relation of belonging
to the same block. A partition is noncrossing if there are no
$a<b<c<d$ with $a,c$ in one block and $b,d$ in a different block.
For two noncrossing partitions $P,Q$ on $[n]$, interleave their vertices as
\[
 1,1',2,2',\ldots,n,n',
\]
putting $P$ on the unprimed positions and $Q$ on the primed positions.
The standard Kreweras complement $K(P)$ is the coarsest $Q$ for which
this combined partition is noncrossing; see Definitions 2.9, 2.17,
and 2.19 of~\cite{kreweras}.

For the other side, an isolated point is a singleton block.
In the standard arc representation, an arc joins consecutive elements
of the same block. A 3-crossing requires three arcs $(a,d),(b,e),(c,f)$
with $a<b<c<d<e<f$; a 3-noncrossing partition has no such configuration.
On two vertices this condition is automatic.

\paragraph{Compute the genuine complement.}
There are exactly two partitions of $\{1,2\}$:
\[
 D=\{\{1\},\{2\}\},\qquad T=\{\{1,2\}\}.
\]
Both are noncrossing. If $P=D$, the primed block $\{1',2'\}$ creates
no crossing, so its coarsest admissible complement is $T$.
If $P=T$, choosing $Q=T$ would give alternating distinct blocks
$\{1,2\}$ and $\{1',2'\}$ in the order $1,1',2,2'$.
Thus $Q=T$ is forbidden, while $Q=D$ is allowed. Consequently
\[
 K(D)=T,\qquad K(T)=D.
\]
Neither partition is fixed. This also shows directly that the map is
an involution in this instance; no general involution assertion is needed.

On the target side, the partition $T$ has no singleton blocks and is
3-noncrossing. It is the only such partition, since the other partition
$D$ has singleton blocks. If the phrase is read as noncrossing partitions
without singletons, or as noncrossing matchings without isolated vertices,
the same two-vertex object $\{1,2\}$ still witnesses nonemptiness.
Thus there cannot even be a surjection from the empty fixed-point set to
the target set. A bijection, with or without a rank condition, is impossible.
\hfill$\square$

\paragraph{Scope.}
We use the standard Kreweras complement and the ordinary meaning of
partitions on $[n]$. The argument uses actual objects at $n=2$, not a
guessed formula for the Riordan numbers. The source imposes no exception
for this value of $n$. The accompanying files are prepared for independent review.

\newpage
\paragraph{Formal verification and exhaustive coverage.}
The accompanying Lean 4.19.0 project imports only \texttt{Std}.
\texttt{Partition n} stores an arbitrary proposition-valued binary
relation on \texttt{Fin n} with proofs of reflexivity, symmetry, and
transitivity. It is therefore an actual set partition, rather than a
label assigned to a preselected example.

\path{every_partition_encoded} proves that every partition of two
points is exactly one of \texttt{model false} and \texttt{model true},
which are $D$ and $T$. \path{model_injective} proves they are distinct,
and \path{every_partition_noncrossing} verifies their eligibility for $NC_2$.
Extensional equality is equality of their full membership relations.

\texttt{interleave} implements the four positions $1,1',2,2'$ using
even and odd indices, with integer division by two giving the original
vertex index. \path{interleave_is_partition} proves directly that the
combined relation is an equivalence relation for any two input partitions.
\texttt{Noncrossing} is the alternating-four-vertex condition above.

\texttt{IsKrewerasComplement P Q} requires both input partitions and
their interleaving to be noncrossing, and requires every admissible
partition to refine $Q$. Its maximality quantifier ranges over
\emph{all} proposition-valued partitions, using the proved complete
classification when finite checking is appropriate.
\path{kreweras_models} verifies the two complements by this definition.
Existence for every input and uniqueness of the complement are proved
separately. Hence \path{no_fixed_points} concerns the fixed points of
the actual uniquely defined operation, not a substitute Boolean map.

The target uses the consecutive-block-element definition of an arc,
the six-vertex 3-crossing condition, and the absence of singleton blocks.
\texttt{targetWitness} constructs $T$ as a valid target object, and
\path{target_unique} proves that every target object is this one.
The theorem \path{conjecture459_refuted} denies the existence of a
surjective function from the actual fixed-point subtype to that target
subtype. This is a necessary condition for the asserted rank-preserving
bijection, so no formalization of an unused rank function is needed.

All finite decisions are kernel checked. There are no admitted proofs,
additional axioms, or native decision procedures. The final theorem
uses only the standard axioms \texttt{propext}, \texttt{Classical.choice},
and \texttt{Quot.sound}. The Lake configuration treats warnings as errors.

\begin{thebibliography}{1}
\bibitem{kreweras}
K. Ebrahimi-Fard and T. Ringeard,
\emph{Planar Binary Trees, Noncrossing Partitions and the Operator-Valued
S-Transform}, Annals of Combinatorics \textbf{29} (2025), 799--836.
Definitions 2.9, 2.17, 2.19.
\url{https://doi.org/10.1007/s00026-024-00730-1}.
\end{thebibliography}
\end{document}
